% =====================================================================
% Problem statement -- carries over the content of the approved topic
% proposal. Please do not change the structure.
% =====================================================================
\vorspannkapitel{Problemstellung}
\abschnittsautor{\alleautoren}

\vorspannabschnitt{Ausgangslage}

Beschreibt hier in wenigen Sätzen, welches Problem im Ausgangszustand besteht,
wen es betrifft und warum eine technische Lösung sinnvoll ist.

\vorspannabschnitt{Untersuchungsanliegen der individuellen Themenstellungen}

\noindent
1) \textbf{Bierochs:} Untersuchung geeigneter Architektur- und Sicherheitskonzepte für ein skalierbares Backend zur Koordination der Datenkommunikation zwischen Applikation, Datenbank und externen Schnittstellen. Erforschung effizienter Verfahren zur serverseitigen Anbindung und Aggregation von Standortdaten sowie Produktdaten-APIs. Zusätzlich Untersuchung technischer Verfahren zur Performance-Optimierung bei der Einbindung von 3D-Grafiken und Ladeanimationen in mobile Applikationen.

\noindent
2) \textbf{Schreiber:} Untersuchung von User interface/User experience-Prinzipien und Interaktionskonzepten zur Entwicklung einer intuitiven, plattformübergreifenden Benutzeroberfläche. Analyse technischer Methoden zur Implementierung eines robusten Kalender- und Erinnerungssystems, insbesondere der Einsatz von Push-Notifications zur rechtzeitigen Benachrichtigung bei Geburtstagen und Anlässen sowie Konzepte zur lokalen Datenspeicherung und Offline-Synchronisation.

\noindent
3) \textbf{Grbic:} Untersuchung geeigneter Datenbankmodelle zur strukturierten Speicherung von Nutzer-, Personen-, Präferenz- und Gesprächsdaten. Vergleich und Evaluierung einer kostenfrei nutzbaren cloudbasierten KI-API mit einem lokal betriebenen Open-Weight-Modell anhand von Kriterien wie Antwortqualität, Latenz, Ressourcenbedarf, Datenschutz und Implementierungsaufwand. Entwicklung geeigneter Prompt-, Kontext- und Ausgabestrukturen für die Generierung personalisierter Geschenkempfehlungen sowie eines anschließenden dialogbasierten Chats zur Anpassung, Erweiterung und Verfeinerung der Vorschläge.


\vorspannabschnitt{Zielsetzung}

\vorspannabschnitt{Geplantes Ergebnis der individuellen Themenstellungen}

\vorspannunterpunkt{1) Bierochs: }
Florian Bierochs: Ein einsatzbereites, abgesichertes Backend-System zur zentralen Datenverarbeitung und Koordination externer Schnittstellen. Es beinhaltet funktionierende API-Anbindungen für lokale Händler sowie an externe Produktdaten-APIs für Geschenkangebote. Zudem werden nahtlos integrierte Ladeanimationen zur gezielten Steigerung der User Experience entwickelt.

\vorspannunterpunkt{2) Schreiber: }
Nicole Schreiber: Ein voll funktionsfähiges, plattformübergreifendes Frontend mit einer modernen und intuitiven Benutzerführung. Die Applikation umfasst ein integriertes Kalender- und Erinnerungssystem für Anlässe und Geburtstage mit funktionierenden Push-Notifications sowie einer zuverlässigen Offline-Synchronisation für die nahtlose Nutzung auf verschiedenen Endgeräten.


\vorspannunterpunkt{3) Grbic: }
Angelo Grbic: Eine strukturierte und performante Datenbankarchitektur für Nutzer-, Profil- und Gesprächsdaten sowie eine evaluierte KI-Empfehlungslogik. Die KI-Anbindung ermöglicht sowohl die initiale Generierung von Geschenkideen als auch einen interaktiven Chat zur Verfeinerung inklusive strukturierter Ausgaben und Filtermöglichkeiten.

\vorspannabschnitt{Meilensteine}

% Deliberately not a table float: the milestones belong directly under the
% heading, not somewhere else in the document.
\vorspannabschnitt{Meilensteine}

% Deliberately not a table float: the milestones belong directly under the
% heading, not somewhere else in the document.
\begin{tabular}{@{}p{0.68\linewidth}p{0.24\linewidth}}
	\toprule
	\bfseries Meilenstein & \bfseries Datum \\
	\midrule
	Einreichung des Diplomantrags & 25.09.2026 \\
	Anforderungsanalyse, Systemkonzept und Gesamtarchitektur abgeschlossen & 28.11.2026 \\
	Funktionstüchtiger MVP mit Frontend-Grundstruktur und Datenbank & 23.01.2027 \\
	Integration/Evaluierung KI-Logik und API-Schnittstellen & 15.02.2027 \\
	Implementierung Push-Notifications, Offline-Sync und Animation-Optimierung & 15.03.2027 \\
	Feature-Complete: vollständige Systemintegration und Refactoring & 15.04.2027 \\
	Systemtests, Qualitätssicherung und finaler Feinschliff & 25.04.2027 \\
	Abgabe der Diplomarbeit (Spätester Abgabetermin) & 03.05.2027 \\
	\bottomrule
\end{tabular}
